\documentclass[12pt,a4paper]{article}
\usepackage[margin=2cm]{geometry}
\usepackage{setspace}
\usepackage{amsmath}
\usepackage{amssymb}
\usepackage{graphicx}

% Set line spacing to 1.5 for readability
\onehalfspacing

\begin{document}

% --- TITLE PAGE ---
\begin{titlepage}
\thispagestyle{empty}
\sffamily

\begin{center}
\vspace*{3.7cm}
{\Large School of Physics, Engineering and Technology}

\vspace{3.1cm}
{\fontsize{24}{28}\selectfont \textbf{BEng Project Initial Report}}

\vspace{1.3cm}
{\fontsize{22}{26}\selectfont \textbf{2025/6}}
\end{center}

\vfill
\noindent\textbf{Student Name:} Facundo Franchino

\vspace{1.2cm}
\noindent\textbf{Project Title:} Structured Riemannian Pruning of Feedback Delay Networks

\vspace{1.2cm}
\noindent\textbf{Supervisors:} Dr. Karolina Prawda

\end{titlepage}

% --- ABSTRACT ---
\begin{abstract}
This thesis proposes a novel, generalisable mathematical and computational method to prune 
feedback delay networks for optimal real-time running (?). The method serves applications
such as;
\item Creating the "fastest white noise generator", that is; pruning an arbitrary 
higher order feedback delay network such as $N=1024$ or $N=4096$.
\item Towards real-time differentiable reverbs. Applications such as being able to create a
real-time viable FDN mimic of a RIR via the combination
of differentiable digital signal processing techniques and the proposed pruning method.
 

\end{abstract}

\clearpage

% --- MAIN CONTENT ---
\section{Introduction}
Artificial reverberation has long been shaped by a 
fundamental trade-off between acoustic realism and computational efficiency,
often described in digital signal processing (DSP) as the quality-efficiency
gap.
At one end of this continuum, convolution with measured
Room Impulse Responses (RIRs) provides a highly faithful representation of a physical space.
However, the need to process long sample sequences results in a
complexity that scales as $O(L \log L)$, where $L$ denotes the impulse-response length.
This incurs significant memory bandwidth and latency penalties, making convolution 
computationally prohibitive for multiple concurrent instances under low-latency real-time 
constraints on consumer-grade central processing units (CPUs).

By contrast, algorithmic reverberators, most notably Feedback Delay Networks (FDNs),
offer computationally efficient recursive structures whose per-sample cost scales
approximately linearly with the number of delay lines, $N$.
Yet traditional low-order FDNs, typically based on $4 \times 4$ or $8 \times 8$
feedback matrices, often exhibit insufficient modal density.
The resulting sparsity of resonant modes gives rise to metallic coloration and audible
ringing artefacts, since the auditory system can resolve individual resonances that do not
blend into a diffuse statistical sound field.






\clearpage



\end{document}